\subsection{深度与投射维数}

定理 1 (Auslander-Buchsbaum).--- 设 A 为 Noether 局部环，M 为具有有限投射维数的有限生成 A-模。有等式

\[
\mathrm{dp} _ {\mathrm{A}} (\mathrm{M}) + \operatorname{prof} _ {\mathrm{A}} (\mathrm{M}) = \operatorname{prof} (\mathrm{A}).
\]

我们对 \(\mathrm{dp}_{\mathrm{A}}(\mathrm{M})\) 作归纳。

\begin{enumerate}
\def\labelenumi{\alph{enumi})}
\item
  若 dp \(_{A}\) (M) 为零，则 M 是非零的有限生成自由 A-模；其深度等于 prof(A) (§ 1, 第 1 节, 注 4)。
\item
  设 \(\mathrm{dp}_{A}(\mathbf{M}) = 1\)，我们选取 M 的一个极小投射分解
\end{enumerate}

\[
0 \to \mathrm{L} _ {1} \xrightarrow {d _ {1}} \mathrm{L} _ {0} \xrightarrow {d _ {0}} \mathrm{M} \to 0
\]

其中 A-模 L₀ 与 L₁ 是非零的有限生成自由模，故深度为 prof(A) (§ 1, 第 1 节, 注 4)。映射 \(1_{\kappa_{A}} \otimes d_{1} : \kappa_{A} \otimes_{A} L_{1} \longrightarrow \kappa_{A} \otimes_{A} L_{0}\) 为零，这蕴含 \(d_{1}\) 属于 \(\mathfrak{m}_{A} \operatorname{Hom}_{A}(L_{1}, L_{0})\)。由 § 1, 第 1 节的注 5，有 \(\operatorname{prof}_{A}(M) = \operatorname{prof}(A) - 1\)。

\begin{enumerate}
\def\labelenumi{\alph{enumi})}
\setcounter{enumi}{2}
\tightlist
\item
  设 \(\mathrm{dp}_{\mathrm{A}}(\mathrm{M}) > 1\)。选取一个正合列
\end{enumerate}

\[
0 \to \mathrm{N} \longrightarrow \mathrm{L} \longrightarrow \mathrm{M} \to 0
\]

其中 L 是有限生成自由 A-模。于是有 \(\operatorname{prof}_{\mathrm{A}}(\mathrm{L}) = \operatorname{prof}(\mathrm{A})\) ( \(\S 1\), 第 1 节, 注 4)，\(\mathrm{dp}_{\mathrm{A}}(\mathrm{N}) = \mathrm{dp}_{\mathrm{A}}(\mathrm{M}) - 1\) (A, X, 第 135 页, 推论 2 c))，故 \(\operatorname{prof}_{\mathrm{A}}(\mathrm{N}) = \operatorname{prof}(\mathrm{A}) - \mathrm{dp}_{\mathrm{A}}(\mathrm{N})\)（归纳假设），特别地 \(\operatorname{prof}_{\mathrm{A}}(\mathrm{N}) < \operatorname{prof}_{\mathrm{A}}(\mathrm{L})\)。于是由 \(\S 1\), 第 1 节的命题 1，得 \(\operatorname{prof}_{\mathrm{A}}(\mathrm{M}) = \operatorname{prof}_{\mathrm{A}}(\mathrm{N}) - 1\)，证明完毕。

注. 鉴于第 3 节命题 4 的推论 2，把定理 1 应用于 A-模 \(\kappa_{\mathrm{A}}\) 就推出：我们处于下述两种情形之一：

\begin{enumerate}
\def\labelenumi{(\roman{enumi})}
\item
  有 \(\mathrm{dp}_{\mathrm{A}}(\kappa_{\mathrm{A}}) = \mathrm{dh}(\mathrm{A}) = +\infty\)
\item
  有 \(\mathrm{dp}_{\mathrm{A}}(\kappa_{\mathrm{A}}) = \mathrm{dh}(\mathrm{A}) = \mathrm{prof}(\mathrm{A}) < +\infty\)。
\end{enumerate}

后文（§ 4, 第 2 节）将看到，(ii) 刻画了正则局部环。

推论 1.--- 保留定理 1 的假设。

\begin{enumerate}
\def\labelenumi{\alph{enumi})}
\item
  我们有 dp \(_{A}\) (M) ≤ prof(A)。为使等号成立，必须且只需极大理想 m \(_{A}\) 与 M 相伴。
\item
  我们有 \(\operatorname{prof}_A(M) \leqslant \operatorname{prof}(A)\)。为使等号成立，必须且只需 M 是自由的。
\item
  事实上，「 \(\operatorname{prof}_{A}(\mathsf{M}) = 0\) 」等价于「 \(\mathfrak{m}_A \in \operatorname{Ass}(A)\) 」(§ 1, 第 1 节, 注 2)。
\item
  事实上，「dp \(_{A}\) (M) = 0」等价于「M 是自由的」。
\end{enumerate}

推论 2. 保留定理 1 的假设，并设 A 还是 Macaulay 环。则 dp \(_{A}\) (M) 是两个正整数 dim(A) - dim \(_{A}\) (M) 与 dim \(_{A}\) (M) - prof(M) 之和。

特别地，这时 \(\mathrm{dp}_{\mathrm{A}}(\mathrm{M})\) 大于 \(\dim (\mathrm{A}) - \dim_{A}(\mathrm{M})\)，且当且仅当 M 是 Macaulay 的时取等号。

推论 3. 设 A 为 Noether 环，M 为具有有限投射维数的有限生成 A-模，i 为一个 \(\geqslant 0\) 的整数，N 为有限生成 A-模，F 为 A-模 \(\operatorname{Ext}_{\mathrm{A}}^{i}(\mathrm{M},\mathrm{N})\)（相应地 \(\operatorname{Tor}_{i}^{A}(\mathrm{M},\mathrm{N})\)）的支撑集。则我们有 \(\operatorname{prof}_{\mathrm{F}}(\mathrm{A}) \geqslant i\)。

事实上，设 \(\mathfrak{p} \in \mathrm{F}\)。由第 2 节的命题 2，有 \(\operatorname{Ext}_{\mathrm{A}_{\mathfrak{p}}}^{i}(\mathrm{M}_{\mathfrak{p}}, \mathrm{N}_{\mathfrak{p}}) \neq 0\)（相应地 \(\operatorname{Tor}_{i}^{\mathrm{A}_{\mathfrak{p}}}(\mathrm{M}_{\mathfrak{p}}, \mathrm{N}_{\mathfrak{p}}) \neq 0\)），故 \(i \leqslant \mathrm{dp}_{\mathrm{A}_{\mathfrak{p}}}(\mathrm{M}_{\mathfrak{p}}) \leqslant \mathrm{dp}_{\mathrm{A}}(\mathrm{M}) < +\infty\) (第 2 节, 命题 3)。定理 1 蕴含 \(\operatorname{prof}(\mathrm{A}_{\mathfrak{p}}) \geqslant i\)。因此 (§ 1, 第 5 节, 命题 8)

\[
\operatorname{prof} _ {F} (A) = \inf _ {\mathfrak {p} \in F} \operatorname{prof} (A _ {\mathfrak {p}}) \geqslant i.
\]

用 § 1, 第 5 节, 注 4 的术语，推论 3 的结论意味着模 \(\operatorname{Ext}_{A}^{i}(M,N)\) 与 \(\operatorname{Tor}_{i}^{A}(M,N)\) 的级 \(\geqslant i\)。它蕴含它们的支撑集在 \(\operatorname{Spec}(A)\) 中的余维数 \(\geqslant i\) ( \(\S 1\) , 第 7 节, 命题 12)。

推论 4.- 设 A 为 Noether Macaulay 环，M 为具有有限投射维数的有限生成 A-模。

\begin{enumerate}
\def\labelenumi{\alph{enumi})}
\tightlist
\item
  设 \(\mathfrak{p} \in \operatorname{Spec}(A)\)。用 \(\mathcal{C}(\mathfrak{p})\) 表示 \(\operatorname{Supp}(M)\) 的包含 \(\mathfrak{p}\) 的不可约分支组成的集合。我们有
\end{enumerate}

\[
\dim_ {A _ {\mathfrak {p}}} \left(M _ {\mathfrak {p}}\right) - \operatorname{prof} _ {A _ {\mathfrak {p}}} \left(M _ {\mathfrak {p}}\right) = \mathrm{dp} _ {A _ {\mathfrak {p}}} \left(M _ {\mathfrak {p}}\right) - \inf _ {X \in \mathscr {C} (\mathfrak {p})} \operatorname{codim} (X, \operatorname{Spec} (A)).
\]

\begin{enumerate}
\def\labelenumi{\alph{enumi})}
\setcounter{enumi}{1}
\item
  映射 \(\mathfrak{p}\mapsto\dim_{\mathrm{A}_{\mathfrak{p}}}(M_{\mathfrak{p}})-\operatorname{prof}_{\mathrm{A}_{\mathfrak{p}}}(M_{\mathfrak{p}})\)（从 Spec(A) 到 \(\overline{Z}\)）是上半连续的。
\item
  \(A\) 中使 A \(_{p}\) -模 M \(_{p}\) 为 Macaulay 的素理想 p 组成的集合在 \(\operatorname{Spec}(A)\) 中是开的且稠密。它与 \(\operatorname{Supp}(M)\) 的交在 \(\operatorname{Supp}(M)\) 中稠密。
\item
  可设 \(\mathfrak{p} \in \operatorname{Supp}(M)\)。令 \(\varphi(\mathfrak{p}) = \dim(A_{\mathfrak{p}}) - \dim_{A_{\mathfrak{p}}}(M_{\mathfrak{p}})\)。由上面的推论 2，我们有
\end{enumerate}

\[
\dim_ {A _ {\mathfrak {p}}} (M _ {\mathfrak {p}}) - \operatorname{prof} _ {A _ {\mathfrak {p}}} (M _ {\mathfrak {p}}) = \mathrm{dp} _ {A _ {\mathfrak {p}}} (M _ {\mathfrak {p}}) - \varphi (\mathfrak {p}).
\]

由于 A 是 Macaulay 环，我们有

\[
\dim (A _ {\mathfrak {p}}) = \operatorname{codim} (V (\mathfrak {p}), \operatorname{Spec} (A)) = \operatorname{codim} (V (\mathfrak {p}), X) + \operatorname{codim} (X, \operatorname{Spec} (A))
\]

对一切 \(X \in \mathcal{C}(\mathfrak{p})\) 成立 (§ 2, 第 2 节, 命题 2 b))；另一方面，我们有 (VIII, § 1, 第 4 节, 命题 9 及第 2 节, 注 3)

\[
\dim_ {A_ {\mathfrak {p}}} \left(\mathrm{M} _ {\mathfrak {p}}\right) = \operatorname{codim} (\mathrm{V} (\mathfrak {p}), \operatorname{Supp} (\mathrm{M})) = \sup _ {\mathrm{X} \in \mathscr {C} (\mathfrak {p})} \operatorname{codim} (\mathrm{V} (\mathfrak {p}), \mathrm{X})
\]

因而

\[
\varphi (\mathfrak {p}) = \inf _ {X \in \mathscr {C} (\mathfrak {p})} \operatorname{codim} (X, \operatorname{Spec} (A)).
\]

\begin{enumerate}
\def\labelenumi{\alph{enumi})}
\setcounter{enumi}{1}
\item
  设 \(\mathfrak{p} \in \operatorname{Spec}(A)\)，\(F\) 为 \(\operatorname{Supp}(M)\) 的不含 \(\mathfrak{p}\) 的诸不可约分支的并。对每个元素 \(\mathfrak{q}\)（属于 \(\operatorname{Spec}(A) - F\)），有 \(\mathcal{C}(\mathfrak{q}) \subset \mathcal{C}(\mathfrak{p})\)，于是由上面的公式得 \(\varphi(\mathfrak{q}) \geqslant \varphi(\mathfrak{p})\)。因此函数 \(\varphi\) 是下半连续的；断言 b) 于是由第 2 节的命题 3 得出。
\item
  设 U 为 Spec(A) 中使 \(M_{p}\) 为 Macaulay 的元素 p 组成的集合。条件 \(p \in U\) 等价于 \(\dim(M_{p}) - \text{prof}(M_{p}) \leqslant 0\)，故由 b) 知 U 是开的。由于 U 包含 \(\text{Spec}(A) - \text{Supp}(M)\)，只需证明 \(U \cap \text{Supp}(M)\) 在 \(\text{Supp}(M)\) 中稠密。对 \(\text{Supp}(M)\) 的每个极小素理想 p，\(A_{p}\) -模 \(M_{p}\) 具有有限长度 (IV, § 2, 第 5 节, 命题 7 的推论 2 及 § 1, 第 3 节, 命题 7 的推论 1)，故为 Macaulay 的；因此 U 与 \(\text{Supp}(M)\) 的每个不可约分支都相交。借助 II, § 4, 第 1 节的命题 1 即可完成证明。
\end{enumerate}

